\begin{helper}[The fair-sign space with a forced escape block]
Let $q\in\mathbb N$. Then there are a measurable space $(\Omega,\mathcal M)$, a
measure $\mu$ on it with $\mu(\Omega)=1$, a filtration
$(\mathcal G_n)_{n\in\mathbb N}$ of sub-$\sigma$-algebras of $\mathcal M$, a
sequence of functions $\xi_n:\Omega\to\R$ and a function
$G:\Omega\to\mathbb N$ such that
\begin{enumerate}
\item $\xi_n(\omega)=1$ or $\xi_n(\omega)=-1$ for every $n\in\mathbb N$ and
  every $\omega\in\Omega$;
\item $\xi_n$ is $\mathcal G_{n+1}$-measurable for every $n\in\mathbb N$;
\item $\mathbb E[\xi_n\mid\mathcal G_n]=0$ almost everywhere, for every
  $n\in\mathbb N$;
\item $\xi_{qG(\omega)+i}(\omega)=1$ for every $\omega\in\Omega$ and every
  $i<q$;
\item $G$ is integrable and $\int_\Omega G\,d\mu\le2^q$.
\end{enumerate}
\end{helper}
\begin{proof}
Write $p=2^{-q}\in(0,1]$.

\emph{Step 1: the product space.}
Let $B=\{\mathrm t,\mathrm f\}$ be a two-point set with its discrete
$\sigma$-algebra and let $\nu$ be the uniform probability measure on $B$, so
that $\nu\{\mathrm t\}=\nu\{\mathrm f\}=\tfrac12$ and $\nu(B)=1$. Let
$\Omega_0=B^{\mathbb N}$ carry the product $\sigma$-algebra and let
$\mu_0=\bigotimes_{k\in\mathbb N}\nu$ be the product measure of the Preamble.
Then $\mu_0$ is a probability measure, and for every finite
$I\subseteq\mathbb N$ and every family of measurable sets
$(T_k)_{k\in I}$ in $B$,
\begin{equation}\label{eq:fse-cyl}
  \mu_0\bigl(\{\omega\in\Omega_0:\omega_k\in T_k\ \text{for all }k\in I\}\bigr)
    =\prod_{k\in I}\nu(T_k).
\end{equation}
In particular the measure of a cylinder prescribing $|I|$ coordinates exactly
is $2^{-|I|}$, and \eqref{eq:fse-cyl} shows that the coordinate maps
$\omega\mapsto\omega_k$ are independent: if $I$ and $J$ are disjoint finite
sets of indices, then every event determined by the coordinates in $I$ is
independent of every event determined by the coordinates in $J$, because both
sides of \eqref{eq:fse-cyl} factorize over $I\cup J$.

\emph{Step 2: the signs on $\Omega_0$.}
Put $Z_0\equiv0$ and, for $k\ge1$,
\[
  Z_k(\omega)=\begin{cases}+1,&\omega_k=\mathrm t,\\-1,&\omega_k=\mathrm f.\end{cases}
\]
Each $Z_k$ is a function of the single coordinate $\omega_k$ and $B$ is
discrete, so each $Z_k$ is measurable, indeed strongly measurable, being
real valued with finite range. Let $\mathcal F_n=\sigma(Z_k:k\le n)$ be the
natural filtration of the sequence $(Z_k)_{k\in\mathbb N}$; since $Z_0$ is
constant, $\mathcal F_n$ is generated by the coordinates $\omega_1,\dots,\omega_n$
and in particular $\mathcal F_n$ is independent of $\sigma(Z_{n+1})$ by Step 1,
$\sigma(Z_{n+1})$ being contained in the $\sigma$-algebra generated by the
single coordinate $\omega_{n+1}$. Moreover, by \eqref{eq:fse-cyl} with
$I=\{k\}$ the law of $\omega_k$ under $\mu_0$ is $\nu$, so for $k\ge1$
\begin{equation}\label{eq:fse-mean}
  \mu_0\{Z_k=1\}=\mu_0\{Z_k=-1\}=\tfrac12,
  \qquad\int_{\Omega_0}Z_k\,d\mu_0=\tfrac12-\tfrac12=0 .
\end{equation}

\emph{Step 3: the escape blocks.}
For $m\in\mathbb N$ let
\[
  A_m=\{\omega\in\Omega_0:\omega_{qm+i+1}=\mathrm t\ \text{for all }i<q\},
  \qquad
  E_m=\bigcap_{j<m}A_j^{\,c} .
\]
Each $A_m$ is a cylinder prescribing the $q$ coordinates with indices in
$(qm,q(m+1)]$, so it is measurable and, by \eqref{eq:fse-cyl},
$\mu_0(A_m)=2^{-q}=p$ and $\mu_0(A_m^{\,c})=1-p$. The set $E_m$ is determined
by the coordinates with indices $\le qm$, because $A_j$ for $j<m$ involves only
coordinates with indices in $(qj,q(j+1)]\subseteq[1,qm]$, while $A_m$ is
determined by the coordinates with indices $>qm$. Hence $E_m$ and $A_m^{\,c}$
are independent by Step 1, and
\[
  \mu_0(E_{m+1})=\mu_0(E_m\cap A_m^{\,c})=\mu_0(E_m)\,(1-p).
\]
Since $E_0=\Omega_0$ has measure $1$, induction on $m$ gives
\begin{equation}\label{eq:fse-geo}
  \mu_0(E_m)=(1-p)^m\qquad(m\in\mathbb N).
\end{equation}

\emph{Step 4: the escape set has full measure.}
Let $S=\bigcup_{m\in\mathbb N}A_m$, a measurable subset of $\Omega_0$. If
$\omega\notin S$ then $\omega\notin A_j$ for every $j$, so $\omega\in E_m$ for
every $m$; thus $\mu_0(S^c)\le(1-p)^m$ for every $m$ by \eqref{eq:fse-geo}.
As $p>0$ we have $0\le1-p<1$, so $(1-p)^m\to0$ and therefore
$\mu_0(S^c)=0$, that is $\mu_0(S)=1$.

Let $\Omega=S$, let $\mathcal M$ be the $\sigma$-algebra induced on $S$ by the
product $\sigma$-algebra, that is the collection of the sets
$\iota^{-1}(B)=B\cap S$ for measurable $B\subseteq\Omega_0$, where
$\iota:S\to\Omega_0$ is the inclusion, and let $\mu$ be the measure induced by
$\mu_0$ on $(\Omega,\mathcal M)$, so that $\mu(\iota^{-1}(B))=\mu_0(\iota(\iota^{-1}(B)))=\mu_0(B\cap S)$
for every measurable $B$; this is legitimate because $S$ is measurable, so
$\iota$ is a measurable embedding and the image of a measurable subset of $S$
is measurable. Since $\mu_0(S^c)=0$ we get the transport identity
\begin{equation}\label{eq:fse-transport}
  \mu(\iota^{-1}(B))=\mu_0(B)\qquad(B\subseteq\Omega_0\ \text{measurable}),
\end{equation}
and in particular $\mu(\Omega)=\mu(\iota^{-1}(\Omega_0))=\mu_0(\Omega_0)=1$.

\emph{Step 5: the signs on $\Omega$ and their conditional means.}
Set $\xi_n=Z_{n+1}\circ\iota$ for $n\in\mathbb N$ and let
$\mathcal G_n=\sigma(Z_{k}\circ\iota:k\le n)$ be the natural filtration on
$\Omega$ of the restricted sequence; as the inverse image of a generated
$\sigma$-algebra is generated by the inverse images,
\begin{equation}\label{eq:fse-filt}
  \mathcal G_n=\{\iota^{-1}(B):B\in\mathcal F_n\}.
\end{equation}
Claim (1) holds by the definition of $Z_k$ for $k\ge1$. Claim (2) holds because
$\xi_n=Z_{n+1}\circ\iota$ is measurable for $\sigma(Z_{n+1}\circ\iota)$, which
is one of the generators of $\mathcal G_{n+1}$ since $n+1\le n+1$.

For claim (3) fix $n$ and note first that, by \eqref{eq:fse-filt},
\eqref{eq:fse-transport} and Step 2, the $\sigma$-algebras
$\sigma(\xi_n)$ and $\mathcal G_n$ are independent under $\mu$: if
$A=\iota^{-1}(A')$ with $A'\in\sigma(Z_{n+1})$ and $B=\iota^{-1}(B')$ with
$B'\in\mathcal F_n$, then
\[
  \mu(A\cap B)=\mu(\iota^{-1}(A'\cap B'))=\mu_0(A'\cap B')
    =\mu_0(A')\,\mu_0(B')=\mu(A)\,\mu(B).
\]
Also $\xi_n$ is bounded and measurable, hence integrable for the probability
measure $\mu$, and by \eqref{eq:fse-transport} and \eqref{eq:fse-mean}
\[
  \mathbb E[\xi_n]=\mu\{\xi_n=1\}-\mu\{\xi_n=-1\}
    =\mu_0\{Z_{n+1}=1\}-\mu_0\{Z_{n+1}=-1\}=0 .
\]
Now let $B\in\mathcal G_n$. Since $\xi_n$ takes only the values $\pm1$,
\[
  \int_B\xi_n\,d\mu=\mu(B\cap\{\xi_n=1\})-\mu(B\cap\{\xi_n=-1\})
    =\mu(B)\bigl(\mu\{\xi_n=1\}-\mu\{\xi_n=-1\}\bigr)=0
    =\int_B0\,d\mu,
\]
where the middle equality used the independence just proved. As the zero
function is $\mathcal G_n$-measurable and integrable, the characterization of
conditional expectation in the Preamble gives
$\mathbb E[\xi_n\mid\mathcal G_n]=0$ almost everywhere, which is claim (3).

\emph{Step 6: the escape index.}
For $\omega\in\Omega=S$ the set $\{m:\omega\in A_m\}$ is nonempty; let
$G(\omega)$ be its least element. If $i<q$ then $\omega\in A_{G(\omega)}$ gives
$\omega_{qG(\omega)+i+1}=\mathrm t$, that is
$\xi_{qG(\omega)+i}(\omega)=Z_{qG(\omega)+i+1}(\omega)=1$, which is claim (4).
By the definition of a least element,
\[
  \{G=m\}=\iota^{-1}(E_m\cap A_m),\qquad
  \{G>m\}=\iota^{-1}(E_{m+1})\qquad(m\in\mathbb N),
\]
so $G$ is measurable, and by \eqref{eq:fse-transport} and \eqref{eq:fse-geo}
\begin{equation}\label{eq:fse-tail}
  \mu\{G>m\}=\mu_0(E_{m+1})=(1-p)^{m+1}\qquad(m\in\mathbb N).
\end{equation}
Since $G\ge0$ is measurable and integer valued, the layer-cake identity of the
Preamble, applied to $G$ and evaluated by splitting
$(0,\infty)=\bigcup_{m\in\mathbb N}(m,m+1]$ on each piece of which
$\lambda\mapsto\mu\{G>\lambda\}=\mu\{G>m\}$ is constant, gives
\[
  \int_\Omega G\,d\mu=\int_0^\infty\mu\{G>\lambda\}\,d\lambda
    =\sum_{m\in\mathbb N}\mu\{G>m\}
    =\sum_{m\in\mathbb N}(1-p)^{m+1}
    =\frac{1-p}{p}\le\frac1p=2^q,
\]
the geometric series being summable because $0\le1-p<1$. In particular the
right-hand side of the layer-cake identity is finite, so $G$ is integrable and
$\int_\Omega G\,d\mu\le2^q$, which is claim (5).
\end{proof}
